% =============================================================
% 4.x การทดสอบระบบ
% =============================================================

\section{การทดสอบระบบ}

\subsection{แผนการทดสอบ}

การทดสอบแบ่งออกเป็น 3 ระดับ ดังตารางที่ \ref{tab:test_plan} โดยในรายงานความก้าวหน้าฉบับนี้ ดำเนินการทดสอบระดับที่ 1 และ 2 แล้ว ส่วนระดับที่ 3 จะดำเนินการกับกลุ่มตัวอย่างจำนวน 10--20 คนในระยะถัดไป

\begin{table}[!htbp]
\centering
\caption{แผนการทดสอบ}
\label{tab:test_plan}
\begin{tabular}{|L{3.10cm}|L{6.11cm}|L{4.07cm}|}
\hline
\textbf{ระดับ} & \textbf{วิธีการ} & \textbf{สถานะ} \\
\hline
1. การทดสอบหน่วย (Unit Testing) & ชุดทดสอบอัตโนมัติด้วย Unity Test Framework (NUnit) ในโหมด EditMode ทดสอบตรรกะที่ไม่ต้องเปิดฉากเกม เช่น ตัวแยกคำสั่ง ระบบช่วยสอน และระบบความก้าวหน้า (ชุด StoryTests ใช้ PlayerPrefs จริง จึงเป็นการทดสอบระดับส่วนประกอบในโหมด EditMode) & ดำเนินการแล้ว 363 กรณี \\ \hline
2. การทดสอบระบบเบื้องต้นโดยผู้จัดทำ (Developer System Testing) & ผู้จัดทำเล่นเกมด้วยไมโครโฟนและการพิมพ์คำสั่ง แล้วตรวจผลจากหน้าจอและแฟ้มบันทึกผล CSV & ดำเนินการแล้วเมื่อวันที่ 8 ตุลาคม 2569 (รอบการเปิดเกมที่มีการสั่งการ 4 รอบ) \\ \hline
3. การทดสอบการยอมรับของผู้ใช้ (UAT) & กลุ่มตัวอย่างตอบแบบวัดความวิตกกังวลก่อนเล่น เล่นเกม แล้วตอบแบบวัดซ้ำและแบบสอบถามประสบการณ์ (ภาคผนวก ก) พร้อมเก็บแฟ้มบันทึกผลแยกตามรอบการเปิดเกม & รอดำเนินการ \\
\hline
\end{tabular}
\end{table}

\subsection{ผลการทดสอบหน่วย}

ชุดทดสอบหน่วยทั้ง 363 กรณี ทดสอบเมื่อวันที่ 9 ตุลาคม 2569 ผ่านทั้งหมด ใช้เวลารวมประมาณ 0.3 วินาที (0.306 วินาที ตามแฟ้มผลการทดสอบของ Unity Test Framework) ดังตารางที่ \ref{tab:unit_result}

\begin{table}[!htbp]
\centering
\caption{ผลการทดสอบหน่วย}
\label{tab:unit_result}
\begin{tabular}{|L{3.50cm}|L{6.42cm}|L{1.36cm}|L{1.56cm}|}
\hline
\textbf{ชุดทดสอบ} & \textbf{สิ่งที่ทดสอบ} & \textbf{จำนวน} & \textbf{ผ่าน} \\
\hline
CommandParserTests & ตัวแยกคำสั่ง: คำสั่งมาตรฐาน องศาหางเสือ คำสั่งผสม การแก้คำสั่งกลางประโยค คำสั่งขัดแย้ง คำถาม และการเรียกชื่อลูกเรือ & 236 & 236 \\ \hline
TeachingTests & การทวนด้วยรูปประโยคมาตรฐาน การอธิบายคำศัพท์ และรายงานภาษา & 36 & 36 \\ \hline
SmallTalkTests & การคุยเล่นกับลูกเรือ และไม่รับคำสั่งเรือเป็นการคุยเล่น & 28 & 28 \\ \hline
StoryTests & โหมดเนื้อเรื่อง 3 บท 9 ภารกิจ ลำดับการปลดล็อก และการบันทึกความก้าวหน้า & 27 & 27 \\ \hline
EmergencyCallTests & การตัดสินคำสั่งฉุกเฉิน เวลาที่ให้ และการนับผล & 14 & 14 \\ \hline
PhraseHintsTests & คำแนะนำประโยคตามสถานการณ์ & 10 & 10 \\ \hline
ProgressionTests & ทักษะภาษา ใบรับรอง และการปลดล็อกอัปเกรด & 9 & 9 \\ \hline
ResearchLogTests & รูปแบบแฟ้ม CSV และการคำนวณเวลาตอบสนอง & 3 & 3 \\
\hline
\multicolumn{2}{|c|}{\textbf{รวม}} & \textbf{363} & \textbf{363 (100\%)} \\
\hline
\end{tabular}
\end{table}

\subsection{ผลการทดสอบตามความต้องการ}

ผลการทดสอบตามความต้องการเชิงฟังก์ชันแสดงดังตารางที่ \ref{tab:test_result}--\ref{tab:test_result2} โดยระบุหลักฐานของแต่ละกรณี กรณีที่ทดสอบได้เพียงบางส่วนจะทดสอบให้ครบในการทดสอบระดับที่ 3

\begin{table}[!htbp]
\centering
\caption{ผลการทดสอบตามความต้องการเชิงฟังก์ชัน (TC-01 ถึง TC-10)}
\label{tab:test_result}
\begin{tabular}{|L{1.12cm}|L{1.21cm}|L{5.04cm}|L{3.36cm}|L{1.68cm}|}
\hline
\textbf{รหัส} & \textbf{FR} & \textbf{รายการทดสอบ} & \textbf{หลักฐาน} & \textbf{ผล} \\
\hline
TC-01 & FR-01, 02 & พูดคำสั่งแล้วได้ข้อความถูกต้อง & การทดลองเล่นของผู้จัดทำ: ถอดเป็นคำสั่งที่ตั้งใจ 2 จาก 7 ครั้ง & ผ่านบางส่วน \\ \hline
TC-02 & FR-03 & ปฏิเสธเสียงที่ไม่มีสัญญาณ ฟังไม่ออก สั้นหรือเบาเกินไป & การทดลองเล่นของผู้จัดทำ (ไม่มีเสียง เสียงรบกวน) และตรวจโค้ด (สั้นหรือเบาเกินไป) & ผ่านบางส่วน \\ \hline
TC-03 & FR-04 & แยกคำสั่งมาตรฐานและคำสั่งผสม & ทดสอบหน่วย 236 กรณี & ผ่าน \\ \hline
TC-04 & FR-05 & ส่งคำสั่งไม่ชัดเจนให้ GPT จำแนกเจตนา & การทดลองวันที่ 11 ต.ค. 2569: GPT-4o mini จำแนกคำสั่งที่ไม่ใช่รูปประโยคมาตรฐาน 2 คำสั่งได้ถูกต้อง & ผ่าน \\ \hline
TC-05 & FR-06 & ทวนคำสั่ง ขอให้พูดซ้ำ และทวนด้วยรูปประโยคมาตรฐาน & ทดสอบหน่วย และการทดลองเล่นของผู้จัดทำ & ผ่าน \\ \hline
TC-06 & FR-07 & สนทนากับลูกเรือโดยอ้างอิงสถานะเรือ & ทดสอบหน่วย และการทดลองวันที่ 11 ต.ค. 2569: ลูกเรือทั้ง 2 คนตอบ 3 คำถามด้วย GPT-4o mini โดยอ้างอิงสถานะเรือ & ผ่าน \\ \hline
TC-07 & FR-08 & แยกลูกเรือที่ถูกเรียกชื่อและแจ้งเมื่อสั่งผิดช่อง & ทดสอบหน่วย (การเรียกชื่อ) & ผ่านบางส่วน \\ \hline
TC-08 & FR-09 & เรือเคลื่อนที่ตามคันโทรเลขและหางเสือ & การทดลองเล่นของผู้จัดทำ & ผ่าน \\ \hline
TC-09 & FR-10 & สุ่มสร้างด่านที่ต่างกันในแต่ละรอบ & การเดินเรืออิสระ 4 รอบระหว่างจับภาพหน้าจอ ได้แผนที่ต่างกันทุกรอบ & ผ่าน \\ \hline
TC-10 & FR-11 & เตือนภัยและตัดสินคำสั่งฉุกเฉิน & ทดสอบหน่วย และการทดลองเล่นของผู้จัดทำ (พิมพ์คำสั่ง ยังไม่ทดสอบด้วยเสียง) & ผ่านบางส่วน \\ \hline
\end{tabular}
\end{table}

\begin{table}[!htbp]
\centering
\caption{ผลการทดสอบตามความต้องการเชิงฟังก์ชัน (TC-11 ถึง TC-19)}
\label{tab:test_result2}
\begin{tabular}{|L{1.12cm}|L{1.21cm}|L{5.04cm}|L{3.36cm}|L{1.68cm}|}
\hline
\textbf{รหัส} & \textbf{FR} & \textbf{รายการทดสอบ} & \textbf{หลักฐาน} & \textbf{ผล} \\
\hline
TC-11 & FR-12 & บทฝึกสอนเริ่มเมื่อเดินเรือครั้งแรก & การทดลองเล่นของผู้จัดทำ (ขั้นตอนแรก) & ผ่านบางส่วน \\ \hline
TC-12 & FR-13 & คำแนะนำประโยคและรายงานภาษา & ทดสอบหน่วย และการทดลองเล่นของผู้จัดทำ & ผ่าน \\ \hline
TC-13 & FR-14 & ใบรับรองและการปลดล็อกอัปเกรด & ทดสอบหน่วย (ตรรกะเท่านั้น ยังไม่ทดสอบหน้าอู่ต่อเรือ) & ผ่านบางส่วน \\ \hline
TC-14 & FR-15 & โหมดเนื้อเรื่องและการบันทึกภารกิจ & ทดสอบหน่วย (ตรรกะเท่านั้น ยังไม่ทดสอบหน้าเลือกภารกิจ) & ผ่านบางส่วน \\ \hline
TC-15 & FR-16 & ขอเทียบท่าและผูกเรือ & ทดสอบหน่วย (แยกคำสั่ง) & ผ่านบางส่วน \\ \hline
TC-16 & FR-17 & แสดงหน้าสรุปผลหลังจบการเดินเรือ & การทดลองเล่นของผู้จัดทำ & ผ่าน \\ \hline
TC-17 & FR-18 & บันทึกแฟ้ม CSV ครบทุกคอลัมน์ & ทดสอบหน่วย และแฟ้มจากการทดลองเล่น พบว่าหัวตารางของ voyages.csv เดิมไม่ปรับตามคอลัมน์ที่เพิ่มขึ้น ต้องแก้ไขและลบแฟ้มเดิมก่อนเก็บข้อมูลจริง & ผ่านบางส่วน \\ \hline
TC-18 & FR-19 & พิมพ์คำสั่งแทนการพูด & การทดลองเล่นของผู้จัดทำ & ผ่าน \\ \hline
TC-19 & FR-20 & บันทึกและเรียกคืนความก้าวหน้า & ทดสอบหน่วย (ตรรกะเท่านั้น ใช้ที่เก็บข้อมูลจำลอง ยังไม่ทดสอบ SettingsService และหน้าจอ) & ผ่านบางส่วน \\
\hline
\multicolumn{4}{|c|}{\textbf{รวม}} & \textbf{ผ่าน 9 ผ่านบางส่วน 10 (จาก 19)} \\
\hline
\end{tabular}
\end{table}

จากตารางที่ \ref{tab:test_result}--\ref{tab:test_result2} กรณีทดสอบผ่าน 9 ผ่านบางส่วน 10 (จาก 19) กรณีที่ผ่านบางส่วน ได้แก่ TC-01, TC-02, TC-07, TC-10, TC-11, TC-13, TC-14, TC-15, TC-17 และ TC-19 โดยมีสาเหตุ 3 กลุ่ม ได้แก่ (1) ความแม่นยำของการถอดเสียงยังต่ำและยังไม่ได้วัดอย่างเป็นระบบ (TC-01) (2) ยังไม่ได้ทดสอบบางเส้นทาง เช่น เสียงสั้นหรือเบาเกินไป และคำสั่งฉุกเฉินด้วยเสียง (TC-02, TC-10) และ (3) ตรรกะภายในผ่านการทดสอบหน่วยแล้ว แต่ยังต้องทดสอบหน้าจอและการทำงานครบทั้งขั้นตอนในเกม (TC-07, TC-11, TC-13, TC-14, TC-15, TC-19) ส่วน TC-17 พบข้อผิดพลาดของหัวตาราง voyages.csv ซึ่งต้องแก้ไขก่อนเก็บข้อมูล กรณีเหล่านี้จะทดสอบให้ครบก่อนและระหว่างการทดสอบระดับที่ 3

สำหรับบริการ OpenAI การทดลองเมื่อวันที่ 11 ตุลาคม 2569 ส่งข้อความพิมพ์ 6 ข้อความให้ลูกเรือระหว่างเดินเรือ พบว่า (1) การสนทนา 3 ข้อความได้คำตอบจาก GPT-4o mini ตามบุคลิกของลูกเรือ ลงท้ายด้วย ``Over'' และอ้างอิงสถานะเรือจริง เช่น ระยะถึงท่าและสภาพทะเล (2) ตัวจำแนกเจตนาแปลง ``Give her everything she has'' เป็นคำสั่งเดินหน้าเต็มกำลังสูงสุด (Flank Speed) ใน 0.84 วินาที และ ``Could you ease back a touch'' เป็นคำสั่งลดความเร็วใน 1.35 วินาที จากนั้นลูกเรือแนะนำรูปประโยคมาตรฐาน ``flank speed'' และ ``reduce speed'' และ (3) คำสั่งมาตรฐาน ``Steady on that course'' ตัวแยกคำสั่งเข้าใจเองโดยไม่เรียกบริการ เวลาของตัวจำแนกเจตนาจึงเพิ่มเข้าไปในเวลาตอบสนองเฉพาะคำสั่งที่ไม่ใช่รูปประโยคมาตรฐาน ทั้งนี้พบว่าลูกเรือตอบว่าเส้นทางข้างหน้าปลอดภัยจากสถานะเรือ ณ ขณะที่ถาม ก่อนที่ระบบเตือนภัยจะแจ้งแนวหินโสโครกในเวลาต่อมา จึงควรชี้แจงผู้เข้าร่วมว่าคำตอบของลูกเรือเป็นข้อมูล ณ ขณะที่ถามเท่านั้น

\subsection{รายละเอียดกรณีทดสอบที่สำคัญ}

\subsubsection{TC-03: แยกคำสั่งมาตรฐานและคำสั่งผสม}

รายละเอียดกรณีทดสอบ TC-03 แสดงดังตารางที่ \ref{tab:tc03}

\begin{table}[!htbp]
\centering
\caption{กรณีทดสอบ TC-03: แยกคำสั่งมาตรฐานและคำสั่งผสม}
\label{tab:tc03}
\begin{tabular}{|L{3.50cm}|L{10.20cm}|}
\hline
\textbf{รายการ} & \textbf{รายละเอียด} \\
\hline
วัตถุประสงค์ & ตรวจว่าตัวแยกคำสั่งเข้าใจคำสั่งมาตรฐานและคำสั่งผสม ใช้คำสั่งสุดท้ายเมื่อผู้เล่นแก้คำสั่งกลางประโยค และไม่ทำคำสั่งเมื่อไม่ควรทำ \\ \hline
ข้อมูลทดสอบ & ``Port twenty'', ``Full ahead and hard to port'', ``Left! No, right!'', ``Full ahead then astern'', ``The tank is full'' \\ \hline
ผลที่คาดหวัง & ได้ค่า 20 องศา; เดินหน้าเต็มกำลังพร้อมหักซ้ายสุด; เลี้ยวขวา (ใช้คำสั่งที่แก้แล้ว); ตัวแยกคำสั่งระบุว่าคำสั่งขัดแย้งและไม่ทำคำสั่ง; ไม่ทำคำสั่งใด เพราะคำว่า full ในประโยคทั่วไปไม่ใช่คำสั่ง \\ \hline
ผลที่ได้ & ตรงตามที่คาดหวังทุกข้อ (ชุดทดสอบ CommandParserTests 236 กรณีผ่านทั้งหมด) \\
\hline
\end{tabular}
\end{table}

\subsubsection{TC-10: เตือนภัยและตัดสินคำสั่งฉุกเฉิน}

รายละเอียดกรณีทดสอบ TC-10 แสดงดังตารางที่ \ref{tab:tc10}

\begin{table}[!htbp]
\centering
\caption{กรณีทดสอบ TC-10: เตือนภัยและตัดสินคำสั่งฉุกเฉิน}
\label{tab:tc10}
\begin{tabular}{|L{3.50cm}|L{10.20cm}|}
\hline
\textbf{รายการ} & \textbf{รายละเอียด} \\
\hline
วัตถุประสงค์ & ตรวจว่าเกมเตือนภัยเมื่อเรือแล่นเข้าหาสิ่งกีดขวาง และตัดสินคำสั่งฉุกเฉินตามความหมาย ไม่ใช่ตามการสะกด \\ \hline
ขั้นตอน & 1) เล่นภารกิจ 1-2 ในโหมดเนื้อเรื่องจนเรือแล่นเข้าหาแนวหินโสโครก 2) ตรวจหน้าจอ 3) ออกคำสั่งหลบ 4) ชุดทดสอบหน่วยส่งคำสั่งหลบทั้งแบบมาตรฐานและแบบไม่มาตรฐาน ทั้งทิศถูกและทิศผิด \\ \hline
ผลที่คาดหวัง & หน้าจอแสดง ``REEF DEAD AHEAD'' พร้อมระยะและทิศทางที่ควรหลบ ยามระวังหน้าเตือนทางวิทยุ และคำแนะนำประโยคเปลี่ยนเป็นคำสั่งหลบ; คำสั่งที่หลบถูกทิศช่วยเรือได้ทั้งแบบมาตรฐานและไม่มาตรฐาน ส่วนคำสั่งที่หลบผิดทิศไม่ช่วยเรือ \\ \hline
ผลที่ได้ & ผ่านบางส่วน ส่วนที่ทดสอบตรงตามที่คาดหวัง ได้แก่ (1) ภาพหน้าจอฉากที่ 4 ในตารางที่ \ref{tab:storyboard2} (2) แฟ้ม emergencies.csv จากการทดลองเล่นของผู้จัดทำ มีสถานการณ์ฉุกเฉิน 2 ครั้ง ผลเป็น saved ทั้ง 2 ครั้ง (ระยะเริ่มต้น 35.92 และ 22.14 เมตร) โดยพิมพ์คำสั่ง ``hard a port'' และ (3) ชุดทดสอบ EmergencyCallTests 14 กรณีผ่านทั้งหมด ส่วนเส้นทางการพูดคำสั่งฉุกเฉินผ่านวิทยุ ยังไม่ทดสอบด้วยเสียง \\
\hline
\end{tabular}
\end{table}

\subsection{ผลการทดสอบความต้องการที่ไม่ใช่เชิงฟังก์ชัน}

เวลาตอบสนองวัดจากการทดลองเล่นของผู้จัดทำเมื่อวันที่ 8 ตุลาคม 2569 ซึ่งมีการกดปุ่มวิทยุ 8 ครั้ง ถอดเสียงได้ 7 ครั้ง (ไม่มีเสียง 1 ครั้ง) ความยาวเสียง 1.42--1.89 วินาที ขณะนั้นเกมยังประมวลผลเสียงด้วยซีพียู พบว่าเวลาแปลงเสียงเป็นข้อความอยู่ระหว่าง 2.47--2.87 วินาที (เฉลี่ย 2.73 วินาที, n = 7) และคำสั่ง ``Stop the engine'' 2 ครั้ง เรือเริ่มทำตามภายใน 2.47 และ 2.56 วินาทีหลังปล่อยปุ่ม ส่วนความแม่นยำของการถอดเสียง ใน 7 ครั้งที่ถอดได้ มีเพียง 2 ครั้งที่ถอดเป็นคำสั่งที่ตั้งใจ (``Stop the engine'' ทั้ง 2 ครั้ง) แต่แฟ้มบันทึกผลไม่มีข้อความที่ผู้เล่นตั้งใจพูดไว้เทียบ และไม่ได้บันทึกเสียง จึงยังคำนวณความแม่นยำไม่ได้ ในการทดสอบระดับที่ 3 ผู้เข้าร่วมแต่ละคนจะอ่านรายการคำสั่งมาตรฐานที่กำหนด 10 คำสั่ง (คำสั่งเป้าหมาย) ระหว่างเล่น และวัดความแม่นยำเป็นร้อยละของคำสั่งเป้าหมายที่ระบบถอดความและแยกคำสั่งได้ถูกต้อง โดยเทียบรายการคำสั่งเป้าหมายกับคอลัมน์ \texttt{heard} และ \texttt{actions} ของแฟ้ม turns.csv ส่วนการเล่นอิสระจะรายงานเพียงร้อยละของการสื่อสารที่ลูกเรือเข้าใจ (\texttt{understood\_pct}) ผลการทดสอบทั้งหมดแสดงดังตารางที่ \ref{tab:nfr_result}

เนื่องจากเวลาแปลงเสียงเป็นส่วนที่ใช้เวลามากที่สุด ผู้จัดทำจึงวัดเวลาแปลงเสียงของ whisper.cpp เปรียบเทียบระหว่างซีพียูกับการ์ดจอบนเครื่องทดสอบในตารางที่ \ref{tab:hardware} เมื่อวันที่ 11 ตุลาคม 2569 โดยใช้เสียงสังเคราะห์ภาษาอังกฤษของ Windows พูดคำสั่งมาตรฐาน 5 คำสั่ง (``Starboard twenty'', ``Hard a port'', ``Slow ahead'', ``Stop the engine'' และ ``Midships'') เติมช่วงเงียบแบบเดียวกับเกม และใช้การตั้งค่าโมเดลเดียวกับเกม ถอดเสียงคำสั่งละ 5 ครั้งโดยไม่นับรอบแรก (Warm-up) ผลแสดงดังตารางที่ \ref{tab:gpu_bench}

\begin{table}[!htbp]
\centering
\caption{เวลาแปลงเสียงเป็นข้อความของ whisper.cpp บนซีพียูและการ์ดจอ (n = 25 ครั้งต่ออุปกรณ์)}
\label{tab:gpu_bench}
\begin{tabular}{|L{4.2cm}|L{2.0cm}|L{2.0cm}|L{2.0cm}|L{2.0cm}|}
\hline
\textbf{อุปกรณ์} & \textbf{ค่ามัธยฐาน} & \textbf{ค่าเฉลี่ย} & \textbf{ต่ำสุด} & \textbf{สูงสุด} \\
\hline
ซีพียู (i5-13420H) & 1{,}603 ms & 1{,}673 ms & 1{,}440 ms & 2{,}254 ms \\ \hline
การ์ดจอ (RTX 4050, Vulkan) & 58 ms & 150 ms & 49 ms & 1{,}935 ms \\
\hline
\end{tabular}
\end{table}

จากตารางที่ \ref{tab:gpu_bench} การ์ดจอลดค่ามัธยฐานของเวลาแปลงเสียงจาก 1.60 วินาทีเหลือ 0.06 วินาที (เร็วขึ้นประมาณ 27 เท่า) และข้อความที่ถอดได้เหมือนกันทุกคำสั่งทั้งสองแบบ ค่าสูงสุดของการ์ดจอเกิดในครั้งแรกของคำสั่งแต่ละความยาว ซึ่งเป็นช่วงที่การ์ดจอเตรียมการคำนวณ ผู้จัดทำจึงเปิดใช้การ์ดจอเป็นค่าเริ่มต้นของเกม ทั้งนี้ผลนี้ใช้เสียงสังเคราะห์และเครื่องเดียว และพบว่าทั้งสองแบบถอดคำสั่ง ``Hard a port'' ผิดเป็น ``Hard Report.'' ซึ่งตัวแยกคำสั่งต้องรองรับคำที่ได้ยินผิดนี้ เวลาตอบสนองจริงของเกมเมื่อใช้การ์ดจอจะวัดซ้ำในการทดสอบระดับที่ 3

\begin{table}[!htbp]
\centering
\caption{ผลการทดสอบความต้องการที่ไม่ใช่เชิงฟังก์ชัน}
\label{tab:nfr_result}
\begin{tabular}{|L{1.45cm}|L{6.85cm}|L{2.80cm}|L{1.74cm}|}
\hline
\textbf{รหัส} & \textbf{ผลการวัด} & \textbf{วิธีตรวจ} & \textbf{ผล} \\
\hline
NFR-01 & เรือเริ่มทำตามคำสั่งภายใน 2.47 และ 2.56 วินาที (เกณฑ์ 3 วินาที) & แฟ้มบันทึกผลการทดลองเล่นของผู้จัดทำ & ผ่านเบื้องต้น (n = 2) \\ \hline
NFR-02 & ไม่มีการส่งเสียงออกนอกเครื่อง และไม่มีการเขียนแฟ้มเสียงของผู้เล่น & ตรวจโค้ด & ผ่าน \\ \hline
NFR-03 & กำหนดเวลารอ 6 และ 15 วินาที และมีเส้นทางขอให้พูดซ้ำเมื่อขัดข้อง เส้นทางไม่มีกุญแจ API ทดสอบแล้วในการทดลองวันที่ 8 และ 11 ต.ค. 2569 (ลูกเรือขอให้พูดซ้ำ) และเมื่อมีกุญแจ API บริการตอบภายในเวลาที่กำหนด (0.84--1.35 วินาที) ส่วนกรณีหมดเวลายังรอทดสอบ & ตรวจโค้ด และการทดลองเล่นของผู้จัดทำ & ผ่านบางส่วน \\ \hline
NFR-04 & กุญแจ API อ่านจากตัวแปรสภาพแวดล้อมหรือแฟ้ม openai.key ที่ระบุใน .gitignore & ตรวจโค้ด & ผ่าน \\ \hline
NFR-05 & ปฏิเสธเสียงที่สั้นหรือเบาเกินไป และตัดเสียงที่ 8 วินาที & ตรวจโค้ด (RadioMic.cs) & รอทดสอบ \\ \hline
NFR-06 & สั่งการด้วยการพิมพ์ได้ (9 ครั้งในการทดลอง) เปลี่ยนปุ่มได้ ปรับขนาดคำบรรยายได้ จอยเกม: ยังไม่ทดสอบ & การทดลองเล่นของผู้จัดทำ & ผ่านบางส่วน \\ \hline
NFR-07 & ทุกประโยคที่แนะนำเป็นรูปประโยคมาตรฐานที่ตัวแยกคำสั่งเข้าใจ & ทดสอบหน่วย & ผ่าน \\
\hline
\end{tabular}
\end{table}

จากตารางที่ \ref{tab:nfr_result} เวลาแปลงเสียงเป็นข้อความคิดเป็นเกือบทั้งหมดของเวลาตอบสนอง ส่วนตัวแยกคำสั่งและการสั่งเรือใช้เวลาเพียงไม่กี่มิลลิวินาที ผู้จัดทำจึงเปิดใช้การ์ดจอประมวลผลเสียงแล้ว (ตารางที่ \ref{tab:gpu_bench}) ผลในตารางที่ \ref{tab:nfr_result}เป็นการวัดเบื้องต้นจากผู้จัดทำคนเดียว (ถอดเสียง n = 7, เวลาเรือตอบสนอง n = 2) ยังสรุปเชิงสถิติไม่ได้ จะวัดซ้ำกับกลุ่มตัวอย่างในการทดสอบระดับที่ 3 ส่วนกรณีหมดเวลาของ NFR-03 ซึ่งต้องจำลองการขัดข้องของเครือข่าย การปฏิเสธเสียงสั้นหรือเบาเกินไปของ NFR-05 และการใช้จอยเกมของ NFR-06 จะทดสอบเพิ่มเติมในระยะถัดไป
